\chapter{Conclusions and future work}
\label{ch:sum}

\section{Conclusion}

This thesis designed and implemented AK\&VEN, an integrated e-commerce and inventory management system for a family
sock business at Dordoi Bazaar. It replaces the manual processes described in the introduction with one application.
Customers browse the real assortment, buy for the whole family, mix any socks to reach a lower price per pair, ask for
a better price in a conversation, and either pay online or reserve their order and pay when they collect it at the
stall or receive it from a courier. The owners and staff manage the catalogue, stock, prices, deliveries, orders and
wholesale requests in an administration interface, and see their sales on a dashboard. The shop and the administration
are one installable web application in three languages; one REST API holds every business rule; and the database
repeats the rules that must never be broken.

The system has two technical contributions. The first is a design in which a language model may talk about money
without being able to decide it. The model only proposes; a deterministic validator limits every proposal to a value
the model never sees, again at checkout; the database refuses an allowed discount larger than its proposal; and the
model's words may not contain a number of their own. Both the proposal and the result are recorded, so the gap between
what the model wanted and what was allowed is visible for every conversation. The second is an inventory that stays
correct and explainable: stock is locked in a fixed order during checkout, reserved stock is held separately from
stock on the shelf, and every change is written to a ledger in the same transaction as the change itself.

\section{Answers to the research questions}

\paragraph{RQ1: Can a language model negotiate while a deterministic component guarantees the limit?} Yes. In the
evaluation no final price went below the limit, whether the assistant was the rule-based one that over-promised six
times or the language model, including during a prompt-injection attempt. The automated tests show the same for an
assistant that deliberately offers 80\%, for offers tampered with at checkout and for rows written directly to the
database. The guarantee rests on code and constraints, not on the model's behaviour.

\paragraph{RQ2: How do the two negotiators compare?} On the sixteen simulated customers the language model was more
cautious: an average final discount of 6\% against 11\%, half of the customers buying against four in five, and the
validator stepping in once against six times. It answered naturally in all three languages within one or two seconds.
The rule-based assistant is therefore the better seller for single messages, and the model the better
conversationalist; its caution comes largely from the prompt, which can be tuned to the owners' wishes.

\paragraph{RQ3: Can the bazaar's habits be expressed in an online shop?} Yes. The shop expresses them as ordinary
interface elements rather than as extra steps: a chat on the product page, a family planner, a price ladder for the
whole bag shown in the bag and on every product, pickup with a code at the container, payment at handover, delivery
areas with clear fees and times, and a case builder with a request inbox for shops. The browser tests and the
accessibility checks show that these work for a first-time visitor, on a phone, and in Russian and Kyrgyz.

\paragraph{RQ4: Can the stock stay consistent and explainable?} Yes. Stock is never read before the affected variants
are locked, always in the same order, so two customers buying the last pair cannot both succeed and two opposite bags
cannot deadlock; this is tested against a real PostgreSQL database over HTTP. Reserved pairs are counted separately,
and a check constraint keeps the reserved quantity from exceeding the stock. Every sale, cancellation, delivery of new
goods and manual correction writes one ledger row in the same transaction, and a test confirms that the rows of each
variant always add up to the stock on the shelf, so any number on the shelf can be traced back to its causes.

\section{Future work}

\begin{itemize}
	\item \textbf{A real card payment provider}, with the payment moved outside the checkout transaction (reserve, pay,
	confirm), and later a payment rail through a Kyrgyz bank. The reservation used for payment at handover already
	provides the first step.
	\item \textbf{A user study} with real customers of the stall, to measure how the two assistants are accepted in full
	conversations rather than single messages.
	\item \textbf{Vector retrieval} for the assistant, using the prepared pgvector table, so it can answer from longer
	product descriptions and the family's notes.
	\item \textbf{A Telegram bot} for the family and for customers. The shop already sends emails at the moments that
	matter, but most of the stall's customers use Telegram and WhatsApp. A bot would send the same messages, let a
	customer check an order by its number, and let the owners mark an order as handed over from their phone.
	\item \textbf{Real photographs} of the products, replacing the generated drawings.
	\item \textbf{Demand forecasting} from the stock ledger, the orders, the waiting lists and the wholesale requests, to
	suggest what and how much to reorder from the suppliers.
	\item \textbf{Translations of the owners' data}, such as section and colour names, into all three languages.
\end{itemize}

\section{Final remarks}

The project started from a simple observation at the stall: nobody pays the label price, and nobody knows exactly
what is left on the shelf. The result is a system that keeps the conversation with the customer, keeps the family in
control of the price, and keeps the stock accurate. It is ready to be used by my parents once the real prices,
photographs and a payment provider are in place.
